\documentclass[10pt,a4paper]{article}
\usepackage[utf8]{inputenc}
\usepackage[T1]{fontenc}
\usepackage{lmodern}
\usepackage[margin=1in]{geometry}
\usepackage{xcolor}
\usepackage{amsmath,amssymb,amsfonts}
\usepackage{graphicx}
\usepackage{booktabs}
\usepackage{adjustbox}
\usepackage{tcolorbox}
\usepackage{enumitem}
\usepackage{microtype}
\usepackage{array}
\usepackage{fancyhdr}

% Custom Colors matching house style
\definecolor{bitsnavy}{RGB}{33,53,94}
\definecolor{bitsblue}{RGB}{44,90,160}
\definecolor{bitsgreen}{RGB}{46,125,82}
\definecolor{bitsamber}{RGB}{200,136,30}
\definecolor{bitsred}{RGB}{178,58,72}

% Callout Boxes
\tcolorboxenvironment{intuition}{colback=bitsblue!8,colframe=bitsblue,title=\textbf{Everyday Picture \& Intuition},fonttitle=\bfseries,boxrule=0.8pt,arc=3pt}
\tcolorboxenvironment{worked}{colback=bitsgreen!8,colframe=bitsgreen,title=\textbf{Fully Worked Example},fonttitle=\bfseries,boxrule=0.8pt,arc=3pt}
\tcolorboxenvironment{everyday}{colback=bitsnavy!5,colframe=bitsnavy,title=\textbf{Real-World Picture},fonttitle=\bfseries,boxrule=0.8pt,arc=3pt}
\tcolorboxenvironment{watchout}{colback=bitsred!8,colframe=bitsred,title=\textbf{Watch Out! Pitfalls},fonttitle=\bfseries,boxrule=0.8pt,arc=3pt}
\tcolorboxenvironment{keytake}{colback=bitsamber!10,colframe=bitsamber,title=\textbf{Key Takeaways},fonttitle=\bfseries,boxrule=0.8pt,arc=3pt}

\newenvironment{intuition}{\begin{tcolorbox}}{\end{tcolorbox}}
\newenvironment{worked}{\begin{tcolorbox}}{\end{tcolorbox}}
\newenvironment{everyday}{\begin{tcolorbox}}{\end{tcolorbox}}
\newenvironment{watchout}{\begin{tcolorbox}}{\end{tcolorbox}}
\newenvironment{keytake}{\begin{tcolorbox}}{\end{tcolorbox}}

\newenvironment{steps}{\begin{enumerate}[label=\textbf{Step \arabic*:},leftmargin=*]}{\end{enumerate}}

\newcommand{\secsub}[1]{\noindent\textit{\large #1}\vspace{0.8em}}
\newcommand{\housefig}[2]{
\begin{figure}[htbp]
\centering
\includegraphics[width=0.85\linewidth]{#1}
\caption{#2}
\end{figure}
}

\pagestyle{fancy}
\fancyhf{}
\rhead{\color{bitsnavy}\bfseries WILP, BITS Pilani}
\lhead{\color{bitsnavy}\bfseries ISM Session 15 Study Companion}
\rfoot{Page \thepage}

\begin{document}

% Banner
\begin{tcolorbox}[colback=bitsnavy,colframe=bitsnavy,arc=0pt,outer arc=0pt,top=12pt,bottom=12pt]
\color{white}
\centering
{\large\bfseries WILP, BITS Pilani}\\[4pt]
{\LARGE\bfseries Session 15: Gaussian Mixture Models \& Practice Problems}\\[4pt]
{\small Course: Introduction to Statistical Methods (ISM) \quad|\quad Instructor: Dr. YVK Ravi Kumar}
\end{tcolorbox}

\vspace{1em}

\section{Gaussian Mixture Models (GMM)}
\secsub{Model complex multimodal data by combining several weighted Gaussian distributions.}

A Gaussian Mixture Model represents an overall population probability density as a linear combination of $K$ underlying Gaussian components:
\begin{equation}
p(x) = \sum_{k=1}^K \pi_k \mathcal{N}(x \mid \mu_k, \Sigma_k)
\end{equation}
where $\pi_k$ are mixing coefficients satisfying $\sum_{k=1}^K \pi_k = 1$ and $\pi_k \ge 0$.

\begin{intuition}
Suppose share prices of Company A average \$90 and Company B average \$100. When combined into one pool, the price distribution shows two distinct peaks. GMM models this combined density as a weighted blend of two bell curves.
\end{intuition}

\housefig{figures/gmm_mixture_density.png}{Bimodal student step count density modeled using a 2-component Gaussian Mixture Model.}

\section{The Expectation-Maximization (EM) Algorithm}
\secsub{Iteratively estimate component parameters and soft cluster responsibilities.}

Because cluster assignments are unobserved latent variables $z_k \in \{0, 1\}$, parameter likelihood cannot be solved in closed form. The Expectation-Maximization (EM) algorithm alternates between two steps until convergence:

\subsection{E-Step (Expectation)}
Compute posterior responsibilities $\gamma(z_k)$ using Bayes' theorem:
\begin{equation}
\gamma(z_k) = P(z_k = 1 \mid x) = \frac{\pi_k \mathcal{N}(x \mid \mu_k, \Sigma_k)}{\sum_{j=1}^K \pi_j \mathcal{N}(x \mid \mu_j, \Sigma_j)}
\end{equation}

\subsection{M-Step (Maximization)}
Re-estimate component parameters using current responsibilities $N_k = \sum_{i=1}^N \gamma(z_{ik})$:
\begin{align}
\mu_k^{\text{new}} &= \frac{1}{N_k} \sum_{i=1}^N \gamma(z_{ik}) x_i \\
\Sigma_k^{\text{new}} &= \frac{1}{N_k} \sum_{i=1}^N \gamma(z_{ik}) (x_i - \mu_k^{\text{new}})(x_i - \mu_k^{\text{new}})^T \\
\pi_k^{\text{new}} &= \frac{N_k}{N}
\end{align}

\housefig{figures/em_responsibilities.png}{Soft cluster responsibilities $\gamma_k(x)$ assigning probabilities across data points.}

\begin{worked}
A wellness app logs step counts (in thousands) of 10 students: $\{3.2, 3.8, 4.1, 4.5, 4.9, 9.2, 9.8, 10.3, 10.7, 11.4\}$.
Component 1 consists of the first 5 values (sedentary); Component 2 consists of the last 5 values (active).
\begin{steps}
\item Calculate mixing coefficients:
\begin{equation}
\pi_1 = \frac{5}{10} = 0.50, \quad \pi_2 = \frac{5}{10} = 0.50
\end{equation}
\item Calculate component means:
\begin{align*}
\mu_1 &= \frac{3.2 + 3.8 + 4.1 + 4.5 + 4.9}{5} = \frac{20.5}{5} = 4.10 \\
\mu_2 &= \frac{9.2 + 9.8 + 10.3 + 10.7 + 11.4}{5} = \frac{51.4}{5} = 10.28
\end{align*}
\item Calculate component variances $s^2 = \frac{\sum (x - \bar{x})^2}{n-1}$:
\begin{align*}
s_1^2 &= \frac{(-0.9)^2 + (-0.3)^2 + 0^2 + 0.4^2 + 0.8^2}{4} = \frac{1.70}{4} = 0.425 \\
s_2^2 &= \frac{(-1.08)^2 + (-0.48)^2 + 0.02^2 + 0.42^2 + 1.12^2}{4} = \frac{2.828}{4} = 0.707
\end{align*}
\item Formulate total GMM PDF:
\begin{equation}
p(x) = 0.50 \mathcal{N}(x \mid 4.10, 0.425) + 0.50 \mathcal{N}(x \mid 10.28, 0.707)
\end{equation}
\end{steps}
\end{worked}

\section{Proportions \& Chi-Square Independence Tests}
\secsub{Test differences in proportions and independence in contingency tables.}

To test proportion differences across samples or test independence in a $2 \times k$ contingency table, we compute expected frequencies $E_{ij} = \frac{R_i C_j}{N}$ and test statistic $\chi^2 = \sum \frac{(O_{ij} - E_{ij})^2}{E_{ij}}$.

\begin{worked}
Employees from Agency 1 (100 total: 67 For, 33 Against) and Agency 2 (150 total: 84 For, 66 Against) were surveyed about a pension plan. Test for equal proportions at $\alpha = 0.01$.
\begin{steps}
\item Overall pooled proportion favoring pension: $\hat{p} = \frac{67 + 84}{100 + 150} = \frac{151}{250} = 0.604$.
\item Expected frequencies: $E_{11} = 100 \times 0.604 = 60.4$, $E_{21} = 150 \times 0.604 = 90.6$.
\item Compute Chi-Square statistic:
\begin{align*}
\chi^2 &= \frac{(67 - 60.4)^2}{60.4} + \frac{(33 - 39.6)^2}{39.6} + \frac{(84 - 90.6)^2}{90.6} + \frac{(66 - 59.4)^2}{59.4} \\
&= 0.721 + 1.100 + 0.481 + 0.733 = 3.035
\end{align*}
\item Critical value at $\alpha = 0.01, df = 1$ is $6.635$. Since $\chi^2 = 3.035 < 6.635$, fail to reject $H_0$. Proportions favoring the plan do not differ significantly between agencies.
\end{steps}
\end{worked}

\begin{watchout}
In GMM parameter estimation, initialization matters! Bad initial guesses for component means can cause the EM algorithm to get stuck in local optima.
\end{watchout}

\begin{keytake}
\begin{itemize}
\item GMM models complex distributions as linear combinations of $K$ Gaussians with mixing weights $\pi_k$.
\item EM algorithm alternates between E-step (computing posterior responsibilities) and M-step (updating means and covariances).
\item Chi-square contingency tests compare observed vs expected frequencies across categorical proportions.
\end{itemize}
\end{keytake}

\end{document}
